% Informe ejecutivo (Entregable 6) — Taller 2, MINE-4101. Objetivo: dos páginas; el detalle está en los notebooks.
% Compilar con XeLaTeX desde esta carpeta:  latexmk   (la configuración está en .latexmkrc)
%
% Toda cifra que se cite sale de la salida de un notebook (la fuente está en un comentario en cada sección).
\documentclass[11pt,a4paper]{article}

% --- Idioma y tipografía -----------------------------------------------------
\usepackage[spanish,es-tabla,es-noshorthands,es-nolists]{babel}
\usepackage{fontspec}
% Fuentes: Red Hat Text/Display si están instaladas; si no, TeX Gyre Heros (incluida en TeX Live).
\IfFontExistsTF{Red Hat Text}{%
  \setmainfont{Red Hat Text}[BoldFont={* SemiBold}, BoldItalicFont={* SemiBold Italic}]%
}{%
  \setmainfont{texgyreheros}[Extension=.otf, UprightFont=*-regular, BoldFont=*-bold,
    ItalicFont=*-italic, BoldItalicFont=*-bolditalic]%
}
\IfFontExistsTF{Red Hat Display}{%
  \newfontfamily\displayfont{Red Hat Display}[BoldFont={* Bold}]%
  \newfontfamily\displaylight{Red Hat Display Light}%
  \newfontfamily\titulofont{Red Hat Display}[BoldFont={* Bold}]%
  \newfontfamily\etiquetafont{Red Hat Display}[BoldFont={* Bold}, LetterSpace=15, WordSpace=2.2]%
}{%
  \newfontfamily\displayfont{texgyreheros}[Extension=.otf, UprightFont=*-regular, BoldFont=*-bold]%
  \newfontfamily\displaylight{texgyreheros}[Extension=.otf, UprightFont=*-regular]%
  \newfontfamily\titulofont{texgyreheros}[Extension=.otf, UprightFont=*-regular, BoldFont=*-bold]%
  \newfontfamily\etiquetafont{texgyreheros}[Extension=.otf, UprightFont=*-regular, BoldFont=*-bold,
    LetterSpace=15, WordSpace=2.2]%
}
\IfFontExistsTF{Adwaita Mono}{\setmonofont{Adwaita Mono}[Scale=0.86]}{}
\usepackage{microtype}
\usepackage[skip=6pt plus 1pt, indent=0pt]{parskip}
\usepackage[autostyle=false,style=english]{csquotes}

% --- Página, color y estructura ----------------------------------------------
\usepackage[a4paper,top=2.2cm,bottom=2.0cm,left=2.2cm,right=2.2cm,headheight=14pt]{geometry}
\usepackage[table]{xcolor}
\definecolor{marino}{HTML}{0B2545}
\definecolor{acento}{HTML}{13A89E}
\definecolor{gris}{HTML}{5B6770}
\definecolor{claro}{HTML}{EEF2F5}
\color{marino!85!black}

\usepackage{graphicx}
\graphicspath{{../figs/}}
\usepackage[section]{placeins}
\makeatletter\setlength{\@fptop}{0pt}\makeatother  % páginas de solo flotantes: alineadas arriba
\usepackage[font=small,labelfont={bf,color=gris},textfont={color=gris},skip=6pt]{caption}
\usepackage{booktabs,tabularx,array,longtable}
\renewcommand{\arraystretch}{1.3}
\arrayrulecolor{marino!60}
\usepackage[inline]{enumitem}
\setlist{itemsep=3pt,topsep=4pt}
\setlist[itemize]{label={\color{acento}\textbullet}}
\setlist[enumerate]{label={\color{acento}\bfseries\arabic*.}}
\usepackage{tikz}
\usepackage{tcolorbox}
\usepackage[framemethod=default]{mdframed}
\usepackage{lastpage}
\usepackage[hidelinks]{hyperref}
\hypersetup{pdftitle={Informe ejecutivo — Predicción de abandono de clientes para FinAlpes},
  pdfauthor={Daniel Triviño, Nicolás Bedoya}}

% --- Títulos -----------------------------------------------------------------
\usepackage{titlesec}
\setcounter{secnumdepth}{0}
\newcommand{\conbarra}[1]{#1\par\nobreak\vspace{3pt}{\color{acento}\rule{2.4cm}{2.2pt}}}
\titleformat{\section}{\displayfont\Large\bfseries\color{marino}}{}{0pt}{}
\titlespacing*{\section}{0pt}{12pt}{2pt}
\titleformat{\subsection}{\displayfont\large\bfseries\color{marino}}{}{0pt}{}
\titlespacing*{\subsection}{0pt}{14pt}{5pt}

% --- Encabezado y pie ----------------------------------------------------------
\usepackage{fancyhdr}
\newcommand{\parteactual}{Informe ejecutivo}
\pagestyle{fancy}
\fancyhf{}
\fancyhead[L]{\small\color{gris}Predicción de abandono de clientes · FinAlpes}
\fancyhead[R]{\small\bfseries\color{acento}\parteactual}
\fancyfoot[L]{\small\color{gris}Taller 2 · MINE-4101 Ciencia de Datos Aplicada}
\fancyfoot[R]{\small\color{gris}\thepage\ / \pageref*{LastPage}}
\renewcommand{\headrule}{\color{claro}\hrule height 0.8pt}
\renewcommand{\footrulewidth}{0pt}

% --- Componentes -------------------------------------------------------------
\usepackage{url}
\DeclareUrlCommand\col{\urlstyle{tt}}                   % nombre de columna del dataset
\def\UrlBreaks{\do\_\do\/\do\.}
\newcommand{\apx}{\ensuremath{\sim}}                     % "aproximadamente"
\newcommand{\etiqueta}[1]{{\etiquetafont\small\bfseries\spaceskip=0.7em\relax\color{acento}\MakeUppercase{#1}}}

% Recuadro con barra lateral: la idea central de una sección.
\newmdenv[topline=false, bottomline=false, rightline=false, linewidth=3pt, linecolor=acento,
  backgroundcolor=claro, innerleftmargin=12pt, innerrightmargin=10pt, innertopmargin=8pt,
  innerbottommargin=8pt, skipabove=6pt, skipbelow=8pt]{destacado}
% Recuadro oscuro: la recomendación principal.
\newtcolorbox{prioridad}{colback=marino, colframe=marino, boxrule=0pt, arc=2pt,
  coltext=white, left=14pt, right=14pt, top=10pt, bottom=10pt}

% Tarjeta de cifra clave: \cifra{número}{qué significa}. Caben tres por línea, separadas con \hfill.
\newcommand{\cifra}[2]{%
  \begin{tcolorbox}[width=0.315\linewidth, height=2.7cm, colback=claro, colframe=acento, arc=0pt,
      toprule=2.5pt, bottomrule=0pt, leftrule=0pt, rightrule=0pt, left=10pt, right=8pt, top=8pt,
      nobeforeafter, box align=top]
    {\displayfont\fontsize{24}{26}\selectfont\bfseries\color{marino}#1}\par\vspace{3pt}
    {\small\raggedright\color{gris}#2\par}
  \end{tcolorbox}}

\newcommand{\figura}[2]{%
  \begin{figure}[htbp]\centering\includegraphics[width=0.92\linewidth]{#1}\caption{#2}\end{figure}}

\begin{document}

% =============================================================================
% Encabezado (en lugar de portada: el informe apunta a dos páginas)
% =============================================================================
\thispagestyle{empty}
\begin{tcolorbox}[colback=marino, colframe=marino, boxrule=0pt, arc=0pt, left=16pt, right=16pt, top=12pt,
    bottom=12pt, coltext=white]
  {\etiquetafont\small\bfseries\spaceskip=0.7em\relax\color{acento}INFORME EJECUTIVO}\par\vspace{6pt}
  {\raggedright\titulofont\fontsize{18}{22}\selectfont\bfseries Predicción de abandono de clientes para FinAlpes}\par
  \vspace{6pt}
  {\small\color{white!80!marino}Daniel Triviño · Nicolás Bedoya \quad|\quad Taller 2 · MINE-4101 · Uniandes \quad|\quad
  9 de noviembre de 2026}
\end{tcolorbox}

% Cifras: 05_evaluacion (PR-AUC 3,2 veces la línea base) y 07_valor (ganancia por cliente y mes de equilibrio).
\begin{destacado}
\textbf{En una frase.} Un modelo que ordena a los clientes por riesgo de abandono permite dirigir la campaña de
retención al \apx16\,\% más riesgoso de la base. En un banco de 20.000 clientes deja \apx COP 271 millones al año y
recupera su costo en menos de tres meses. Recomendamos implementarlo primero como piloto.
\end{destacado}

\cifra{3,2$\times$}{mejor que contactar al azar para encontrar a quienes se van (PR-AUC)}\hfill
\cifra{13.500}{COP de ganancia esperada por cliente evaluado, al año}\hfill
\cifra{2,7 meses}{para recuperar la inversión de \apx COP 61 millones}

% =============================================================================
\section{El problema y los datos}
% Fuentes: Taller2.pdf, 01_entendimiento, 02_limpieza.

FinAlpes quiere detectar a tiempo a los clientes que van a abandonar el banco para ofrecerles una campaña de
retención. La campaña convence al 30\,\% de quienes iban a irse, pero se paga por cada cliente contactado. Contactar a
todos es caro, y contactar al azar desperdicia la mayoría de las ofertas. La pregunta es \textbf{a quién contactar y
si vale la pena}.

Usamos una muestra de 10.045 clientes europeos de un banco, con 11 variables de perfil comercial. El 20,4\,\%
abandonó. La muestra traía errores: el 11\,\% de las filas tenía algún dato faltante, y había 45 filas duplicadas, 95
países o géneros mal escritos y 15 filas con valores imposibles (por ejemplo, edades de 130 años). Tras la limpieza
quedaron 9.985 clientes. Los faltantes se completan dentro del modelo.

% =============================================================================
\section{El modelo}
% Fuentes: 04_modelado, 05_evaluacion.

Comparamos tres modelos (regresión logística, SVM y XGBoost) con el 80\,\% de los datos, y reservamos el 20\,\%
restante para una única evaluación final. Los elegimos por su capacidad de \emph{ordenar} a los clientes por riesgo
(PR-AUC). La exactitud no sirve aquí: decir que nadie abandona ya acierta el 79,6\,\%. \textbf{XGBoost fue el mejor}:
PR-AUC de 0,662 en prueba, 3,2 veces la de contactar al azar. Por decisión ética, el modelo \textbf{no usa el género
ni el país} del cliente.

% =============================================================================
\section{Hallazgos}
% Fuentes: 05_evaluacion (umbral 0,33, lift, auditoría), 06_interpretabilidad (SHAP, ICE).

\begin{enumerate}
  \item \textbf{El modelo concentra el abandono.} En la base de prueba, 6 de cada 10 clientes que marca sí iban a
    irse, frente a 2 de cada 10 al azar. Además, alcanza a más de la mitad de quienes se van. El 10\,\% más
    riesgoso concentra el 38\,\% de los abandonos.
  \item \textbf{Tres variables explican casi todo.} La edad, el número de productos y la actividad del cliente
    suman el 90\,\% del peso de las predicciones. El riesgo sube con la edad hasta los 50--59 años. Con 2 productos
    es mínimo, y con 3 o 4 se dispara.
  \item \textbf{La inactividad es la señal más accionable.} Un cliente inactivo de 55 años tiene, en promedio, un
    riesgo del 68\,\%, frente al 30\,\% de uno activo de la misma edad. Es una asociación, no una causa: activar al
    cliente no garantiza retenerlo.
  \item \textbf{Hay puntos ciegos.} El modelo casi no detecta el abandono de clientes con 2 productos o menores de 40
    años, y subestima el riesgo en Alemania y en mujeres. Para verlos hacen falta datos de comportamiento
    (transacciones, reclamos, uso de canales), no otro algoritmo.
\end{enumerate}

% =============================================================================
\section{Cuánto vale para el banco}
% Fuente: 07_valor (Entregable 5). Método: diapositivas de generación de valor (Provost y Fawcett, cap. 7) y
% Práctica 4, act. 3 del sitio del curso.

Cada acierto vale $r \cdot V - c$, donde $r$ es la tasa de éxito de la oferta, $V$ el ingreso anual del cliente y
$c$ el costo de la oferta: se retiene al 30\,\% y se paga la oferta. Cada falsa alarma cuesta $c$. Los supuestos
para un banco promedio son:

{\small\renewcommand{\arraystretch}{1.15}
\begin{tabularx}{\linewidth}{@{}>{\raggedright}p{4.4cm} >{\raggedright}p{4.3cm} >{\raggedright\arraybackslash}X@{}}
\toprule
\textbf{Supuesto} & \textbf{Valor} & \textbf{Fuente} \\
\midrule
Clientes · abandono anual & 20.000 · 15\,\% & Enunciado (promedio de la banca minorista) \\
Éxito de la oferta ($r$) & 30\,\% de los que iban a irse & Enunciado \\
Ingreso anual por cliente ($V$) & COP 500.000--900.000 (central: 700.000) & Enunciado (reporte 6-K de Nu, 2026) \\
Costo por contacto ($c$) & $V/30$ \apx COP 23.300 & Ejercicio de churn del curso (oferta de 50 por cliente de
  1.500) \\
Inversión inicial & \apx COP 61 millones & Presupuesto de la Práctica 4 del curso, con el costo-hora de un
  científico de datos colombiano \\
\bottomrule
\end{tabularx}}

\begin{itemize}
  \item \textbf{Ganancia esperada:} \apx COP 13.500 por cliente evaluado y por año, \apx COP 22,6 millones al mes en
    toda la base. La inversión se recupera en \textbf{2,7 meses}, y el ROI a 12 meses es del 344\,\%.
  \item \textbf{Margen de seguridad:} cada contacto es rentable mientras más del 11\,\% de los contactados iban a
    irse. Con la tasa de abandono del banco, el modelo acierta en \apx1 de cada 2.
  \item \textbf{Sensibilidad:} la ganancia es positiva en 11 de las 12 combinaciones de ingreso (COP 500.000--900.000)
    y costo de la oferta (COP 10.000--100.000). Solo se pierde dinero si la oferta cuesta COP 100.000 y el cliente
    deja COP 500.000. El supuesto que más pesa es el \textbf{costo de la oferta}.
  \item \textbf{Umbral:} contactar al 40\,\% de la base en lugar del 16\,\% sube la ganancia \apx18\,\%, si el
    banco tiene capacidad para esa campaña.
\end{itemize}

% =============================================================================
\section{Recomendación}

\begin{prioridad}
\textbf{Implementar el modelo como piloto.} Puntuar la base una vez al año y contactar a los clientes con riesgo
mayor o igual a 0,33 (\apx16\,\% de la base, unos 3.100 clientes), reservando un grupo de control sin oferta. El
piloto mide lo que hoy es un supuesto: el costo real de la oferta y cuántos clientes retiene. Antes de escalar, hay que
reentrenar el modelo con clientes del propio banco y sumar datos de comportamiento.
\end{prioridad}

% =============================================================================
\section{Limitaciones}

\begin{itemize}
  \item \textbf{Otro mercado.} El modelo aprendió de clientes europeos en un solo momento. No sabemos si acierta
    igual con clientes colombianos ni con el paso del tiempo.
  \item \textbf{Supuestos sin validar.} El costo de la oferta y la inversión vienen de ejemplos del curso, sin
    costos recurrentes. Se usó el ingreso y no el margen, con un solo año de permanencia y el mismo 30\,\% de
    éxito para todos.
  \item \textbf{Asociación, no causalidad.} El modelo dice a quién contactar, no por qué se va el cliente ni si la
    oferta lo hará quedarse.
\end{itemize}

{\small\color{gris}El detalle está en los notebooks \col{01_entendimiento} a \col{07_valor} del repositorio. La
declaración de uso de IA generativa (Entregable 7) está en \col{declaracion_IA.md}.}

\end{document}
